\chapter*{Preface}
% ##################################################################################################################

\begin{center} \includegraphics[width=0.25\textwidth, angle=0]{figures/MATSimBook.png} \end{center}

% ##################################################################################################################
The intense MATSim development and research during the past 10 years has generated a large body of knowledge composed of a vast amount of studies, dissertations, and projects in research and practice. Time to put together and contextualize all the pieces in a summary book.

There have been \ah{33} authors and 3 editors contributing to this book. Chapters have been written by the respective software contributor whenever possible, which should ensure capturing the most complete and detailed knowledge. 

The book is intended to give the new MATSim user a quick start in running MATSim. It shall furthermore provide the experienced MATSim user and the MATSim developer with details on how to extend MATSim by plugging in the available modules and by programming against the MATSim API. Last, but not least, the book wants to contribute to the methodological discussion on the relatively novel field of microsimulation by compiling our conceptual insights on MATSim gained over the years.

Aiming for the purpose defined above, the book is divided into the three parts focused on \emph{using, extending} and \emph{understanding} MATSim.

\textbf{Part I: using MATSim} ... \\
Chapter~\ref{ch:introducing} introduces the MATSim basics including its traffic flow model and the underlying co-evolutionary principle. Chapter~\ref{ch:lgstarted} takes the MATSim novice at hand and shows him how to set up and run a basic MATSim scenario. Scoring is a very central part in MATSim; a full chapter, Chapter~\ref{ch:scoring}, is thus assigned to scoring. At the point when the reader is ready to set up his own real-world MATSim scenario, Chapter~\ref{ch:scenarios} shows him or her the scenarios that have been implemented at other places worldwide. 

% -----------------
\textbf{Part II: extending MATSim} ... \\
Chapter~\ref{ch:modules} introduces MATSim modules and explains how to use the available modules introduced in Chapters~\ref{ch:pt} through \ref{ch:businessanalytics}. Chapter~\ref{ch:developmentprocess} shortly describes how new functionality usually enters MATSim. Chapter~\ref{ch:extensionpoints} goes one step further and explains how to write your own MATSim module and how to contribute it to MATSim. It details at which points MATSim can be extended, and it digs a little deeper into MATSim and provides details about the very central MATSim concept of ``events''. How to add listeners and use a customized controller is also explained here.

% -----------------
\textbf{Part III: understanding MATSim} ... \\
Chapter~\ref{ch:basicprocedure} formulates the basic MATSim procedure, in more mathematical terms, as a fixed point problem. Chapter~\ref{ch:systemspec} describes MATSim's basic system characteristics, including variability of simulation results and aggregation levels, and it localizes it in in the spectrum of traffic analysis tools. Chapter~\ref{ch:abta} elaborates on agent-based traffic assignment, and Chapter~\ref{ch:montecarlo} continues on that by presenting MATSim as a fundamentally stochastic tool based on random distributions and suitably describable as a Monte Carlo engine. Chapter~\ref{ch:discretechoice} investigates MATSim within the framework of discrete choice modeling, and Chapter~\ref{ch:eco} provides an economic view on MATSim. Chapter~\ref{ch:kinematicwaves} analyzes MATSim's traffic flow model in relation to kinematic waves.

The book is concluded by the presentation of promising research avenues (Chapter~\ref{ch:researchavenues}).

% ##################################################################################################################
\section*{Related Material}
The book is focused on the more stable aspects of MATSim application and development, while the practical issues, such as specific configurations or code details---usually being part of manuals---can be found in the MATSim user guide (shipped with the MATSim releases) and provided on the MATSim web page \citep[][]{MATSim_Userguide_2014}.

% ##################################################################################################################
\section*{Acknowledgments}
editors -> authors \\
editors + authors -> copy editor \\
-> Prof. Axhausen for financing of 1 day/week for Horni.editor\\
-> SNF for publication grant (if yes) \\
-> all developers of recent years and lead of Prof. Nagel and Axhausen

% ##################################################################################################################
